\capitulo{1}{Introducción}

La monitorización de ecosistemas a partir de imágenes aéreas se ha convertido en una herramienta clave para detectar dinámicas espaciales asociadas a la degradación ambiental y a la pérdida de biodiversidad~\cite{turner2003remoteBiodiversity, pettorelli2014appliedEcologists, cavenderbares2022integratingRemoteSensing}. En este contexto, las sendas de herbívoros, originadas por el pisoteo continuado y organizadas en redes de trazado complejo, constituyen un indicador relevante de actividad de herbivoría intensa y, por tanto, un potencial marcador de zonas de presión ecológica~\cite{diezpastora2025trails}. Sin embargo, su identificación manual en ortoimágenes resulta costosa, subjetiva y difícil de escalar~\cite{pettorelli2014appliedEcologists, diezpastora2025trails}, mientras que aproximaciones automáticas clásicas tienden a depender de información topográfica adicional o a producir predicciones a nivel de parche, lo cual limita la delimitación precisa de la estructura y continuidad de las sendas~\cite{diezpastora2025trails}. Frente a ello, la segmentación semántica permite abordar el problema como una clasificación densa a nivel de píxel, incorporando contexto local y global para reconstruir de forma más fiel la geometría de estos patrones sobre imágenes RGB~\cite{long2015fcn, minaee2022segmentationSurvey}.

Este Trabajo Fin de Grado parte del algoritmo descrito en el artículo base~\cite{diezpastora2025trails}, que evalúa arquitecturas modernas de segmentación semántica para mapear sendas de herbívoros en imágenes aéreas y pone de manifiesto la viabilidad de detectar estas estructuras mediante modelos de aprendizaje profundo. A partir de dicha base, el objetivo principal del proyecto consiste en trasladar esta capacidad de segmentación a un entorno accesible para el usuario final. Esto será posible mediante el desarrollo de una aplicación web que permita cargar una imagen, ejecutar el proceso de inferencia y devolver una visualización interpretables, superponiendo y dibujando las sendas detectadas sobre la imagen original. La aplicación incorporará un flujo de uso orientado a la interacción iterativa, de modo que el usuario pueda eliminar la imagen procesada y repetir el procedimiento con nuevas entradas, consolidando así un prototipo funcional que encapsule el \textit{pipeline} completo, desde la recepción del dato hasta la presentación del resultado.

Adicionalmente, el trabajo se plantea con una proyección evolutiva hacia un escenario de mayor escalabilidad, donde la fuente de datos no sea únicamente una imagen aislada subida por el usuario, sino servicios cartográficos de alta resolución, como un entorno soportado por Google Earth Engine~\cite{gorelick2017google, googleearthengine_web} o por recursos del Plan Nacional de Ortofotografía Aérea (PNOA)~\cite{pnoa_ign}. En esa extensión, el usuario seleccionará una zona de interés sobre el terreno y el sistema obtendría la imagen correspondiente para segmentarla y representar automáticamente la red de sendas en el área escogida. Esto abre la puerta a funcionalidades complementarias como la descarga del resultado, el envío por correo electrónico o su integración en flujos de trabajo de análisis territorial.

La memoria se organiza de la siguiente manera: en primer lugar, se presenta el marco teórico necesario para contextualizar la segmentación semántica y los elementos arquitectónicos que la sustentan; a continuación, se describen las arquitecturas consideradas y el enfoque metodológico adoptado para su utilización en el problema de detección de sendas; seguidamente, se detallan el diseño e implementación del sistema, incluyendo el \textit{pipeline} de procesamiento, la construcción de la interfaz web y la estrategia de despliegue; posteriormente, se exponen las pruebas realizadas y los resultados obtenidos, junto con su análisis; finalmente, se recogen las conclusiones y se definen líneas de trabajo futuro orientadas a la explotación de fuentes geoespaciales y a la ampliación de funcionalidades. Como materiales entregables complementarios a la memoria, se incluyen el código fuente del prototipo junto con su correspondiente \textit{mockup}, los recursos necesarios para su ejecución y una guía de uso y reproducción del entorno, con el fin de facilitar la evaluación, la mantenibilidad y la continuidad del proyecto.

